\subsection{A linear depth bound with zero biases}

\begin{theorem}[Unrestricted zero-bias critical networks]
\label{thm:zero-bias-linear}
For any zero-bias network with absolute row sums at most one, there
is an exact sampler with the preprocessing bounds of
Theorem~\ref{thm:upper} and
\[
 \E N_{\rm calls}\le128(D+1),\qquad
 \E[\hbox{online bit work}]=O(\bits^4(D+1)).
\]
Weights can have arbitrary signs. The bound is uniform in the context.
\end{theorem}
\paragraph{Proof idea.}
With zero biases, every stored center is zero and no bias race is needed.
The output-radius gate and the majority costs then telescope to a linear
bound in depth.

\begin{proof}
Use reference gain one and the stored radii of
\eqref{eq:rounded}. Store every center as exactly zero, which is
valid because all biases are zero and tanh is odd. At a row of
radius $\rho$, the normalized call has an active gate
$\tanh\rho/r_\ell\le1$, followed by the majority factory of
Lemma~\ref{lem:arity}; there is no bias race. Its expected number
of reached children is at most
\[
 1+\rho^2+\rho^4\le\exp(r_{\ell-1}^2+r_{\ell-1}^4).
\]
The actual top output retains its gate $\tanh\rho\le r_D$.

Write $q_j=3/(2j+3)$. Lemma~\ref{lem:rounding} gives
$0\le r_j^2-R_j^2\le6D\varepsilon$ and $r_j\le1$.
The derivative of $v+v^2$ on $[0,1]$ is at most three, so the
rounding contribution to any suffix logarithm is at most
$18D^2\varepsilon<1/50$.
The integral estimates in \eqref{eq:qsums} give, for $0\le k<D$,
\[
 \sum_{j=k}^{D-1}(r_j^2+r_j^4)
 \le \frac32\log\frac{D+1}{k+1}+\frac9{4(k+1)}+\frac1{50}
 \le \frac32\log\frac{D+1}{k+1}+3.
\]
The conditional charging lemma therefore bounds the expected total
number of calls, with the top call charged separately, by
\[
 1+r_D\sum_{k=0}^{D-1}
       e^3\left(\frac{D+1}{k+1}\right)^{3/2}
 \le1+2e^3(D+1)\sum_{k\ge1}k^{-3/2}
 \le128(D+1).
\]
Here $r_D\le2/\sqrt{D+1}$, $\sum k^{-3/2}\le3$, and $e^3<21$.
The scalar and index implementations in Section~\ref{sec:bits}
cost $O(\bits^4)$ conditionally per reached call. Their geometric
refinement tails and Lemma~\ref{lem:stopping} give the bit bound.
The majority identity proves exactness; finite depth and almost
sure termination of each scalar call prove almost sure termination
of the entire sampler.
\end{proof}
